\abstracten

With the rapid advancement of Large Language Models (LLMs), automated workflow systems centered on multi-agent collaboration have demonstrated broad application prospects in complex domain-specific tasks. However, existing workflow generation methods face two key challenges in professional domains such as power systems. First, generation often relies excessively on the parametric knowledge of LLMs and fails to exploit structured knowledge from domain documents, equipment relations, and operating procedures. Second, the mapping from subtasks to agents lacks quantitative modeling of capabilities, tools, and interfaces, which may lead to inaccurate role assignment and limited workflow interpretability.

To address these challenges, this thesis proposes an autonomous agent workflow generation method that combines Graph Retrieval-Augmented Generation (GraphRAG) with three-dimensional quantitative subtask-agent matching. A multimodal knowledge graph construction pipeline is designed to integrate power-domain texts, images, and flowcharts into a unified graph structure through entity-relation extraction, vision-language descriptions, and BGE-M3 vector embeddings, with the graph stored in Neo4j. A subgraph-oriented GraphRAG retrieval method is then developed, where semantic vector search retrieves seed nodes and breadth-first traversal supplements multi-hop neighborhood relations to form structured graph context. Furthermore, a hybrid subtask-agent matching algorithm is constructed with a weighted scoring function over capability embedding similarity, tool coverage, and input-output compatibility, enabling explainable mapping from decomposed subtasks to specialized agents. Finally, a workflow Directed Acyclic Graph (DAG) orchestration and validation mechanism with iterative repair is designed to support topological sorting, cycle detection, interface conformance, and LangGraph export.

Experiments are conducted on a stratified benchmark of 30 tasks covering grid fault diagnosis (OPS), power knowledge question answering (KQA), and power data analysis (DA), using PowerGridQA and the OPSD open power time-series dataset. The results show that: (1) all methods achieve 100\% workflow executability, verifying the effectiveness of the three-layer DAG validation; (2) the three-dimensional matching algorithm is the decisive factor of role assignment accuracy (RAA), improving the average RAA by 36.9 percentage points over the no-matching baseline (0.056 $\rightarrow$ 0.425), with consistent gains across all three scenarios; (3) the BGE-M3 semantic upgrade of the step completeness metric raises measured coverage from the 0.18 level (lexical version) to the 0.77 level, more faithfully reflecting paraphrased Chinese power-domain steps; (4) breadth-first subgraph expansion does not yield stable gains on the current 4{,}337-node power knowledge graph and even introduces noise nodes that perturb downstream task decomposition, suggesting that adaptive hop control conditioned on query confidence is a promising direction. The findings verify the effectiveness of graph-based semantic retrieval combined with quantitative agent matching for workflow generation in professional domains.

\vspace{5pt}

\keywordsen Multimodal Knowledge Graph, Graph Retrieval-Augmented Generation, Agent Workflow, Subtask-Agent Matching, Power Systems
